\vspace{0.5em}\noindent\textbf{AWS Networking Advanced: 3 kỹ thuật nâng cấp VPC chuẩn doanh
nghiệp}\par

\vspace{0.5em}\noindent\textbf{Giới thiệu}\par

Mô hình Amazon VPC cơ bản gồm Public Subnet, Private Subnet và NAT
Gateway là nền tảng phù hợp để bắt đầu triển khai ứng dụng trên AWS.
Kiến trúc này giúp tách các tài nguyên cần truy cập trực tiếp từ
Internet khỏi những thành phần nội bộ như application server, database
và cache.

Tuy nhiên, khi hệ thống mở rộng thành nhiều môi trường như Development,
Staging và Production hoặc được chia thành nhiều microservices và VPC
khác nhau, kiến trúc cơ bản có thể bắt đầu xuất hiện ba hạn chế lớn:

\begin{itemize}
\tightlist
\item
  Chi phí NAT Gateway tăng cao khi các tài nguyên trong Private Subnet
  thường xuyên truy cập Amazon S3, Amazon DynamoDB và những dịch vụ AWS
  khác.
\item
  Rủi ro Single Point of Failure nếu toàn bộ lưu lượng outbound chỉ phụ
  thuộc vào một NAT Gateway trong một Availability Zone.
\item
  Việc quản lý kết nối trở nên phức tạp khi nhiều VPC được kết nối với
  nhau bằng VPC Peering theo mô hình mạng lưới.
\end{itemize}

Bài viết này trình bày ba kỹ thuật nâng cấp kiến trúc VPC theo hướng phù
hợp hơn với hệ thống doanh nghiệp, bao gồm VPC Endpoints, Multi-AZ
networking và AWS Transit Gateway.

\vspace{0.5em}\noindent\textbf{1. Giảm chi phí và tăng bảo mật với VPC
Endpoints}\par

\vspace{0.5em}\noindent\textbf{Vấn đề khi sử dụng NAT
Gateway}\par

Các EC2 instances, container hoặc application servers nằm trong Private
Subnet thường không có địa chỉ IP công cộng. Khi những tài nguyên này
cần truy cập Amazon S3, DynamoDB hoặc các AWS services khác, lưu lượng
có thể phải đi qua NAT Gateway.

Luồng kết nối cơ bản có thể được mô tả như sau:

\begin{Shaded}
\begin{Highlighting}[]
\NormalTok{Private Subnet}
\NormalTok{      |}
\NormalTok{      v}
\NormalTok{NAT Gateway}
\NormalTok{      |}
\NormalTok{      v}
\NormalTok{AWS Public Service Endpoint}
\end{Highlighting}
\end{Shaded}

Cách triển khai này có thể phát sinh phí NAT Gateway theo giờ và phí xử
lý dữ liệu. Lưu lượng cũng phải đi qua public service endpoint dù cả ứng
dụng và dịch vụ đích đều hoạt động bên trong hạ tầng AWS.

\vspace{0.5em}\noindent\textbf{Gateway Endpoint}\par

Gateway Endpoint được sử dụng cho hai dịch vụ chính:

\begin{itemize}
\tightlist
\item
  Amazon S3.
\item
  Amazon DynamoDB.
\end{itemize}

Sau khi Gateway Endpoint được tạo và liên kết với route table của
Private Subnet, lưu lượng đến S3 hoặc DynamoDB có thể được chuyển trực
tiếp qua mạng nội bộ của AWS mà không cần đi qua NAT Gateway.

Luồng kết nối mới:

\begin{Shaded}
\begin{Highlighting}[]
\NormalTok{Private Subnet}
\NormalTok{      |}
\NormalTok{      v}
\NormalTok{VPC Gateway Endpoint}
\NormalTok{      |}
\NormalTok{      v}
\NormalTok{Amazon S3 hoặc Amazon DynamoDB}
\end{Highlighting}
\end{Shaded}

Gateway Endpoint mang lại các lợi ích:

\begin{itemize}
\tightlist
\item
  Không cần sử dụng NAT Gateway cho lưu lượng đến S3 và DynamoDB.
\item
  Giữ lưu lượng bên trong mạng AWS.
\item
  Giảm chi phí xử lý dữ liệu qua NAT Gateway.
\item
  Có thể kiểm soát quyền truy cập bằng endpoint policy.
\item
  Giảm việc phụ thuộc vào kết nối Internet outbound.
\end{itemize}

Gateway Endpoint cho S3 và DynamoDB không thu phí theo giờ sử dụng
endpoint, vì vậy đây thường là một trong những giải pháp đầu tiên nên
cân nhắc khi tối ưu chi phí network.

\vspace{0.5em}\noindent\textbf{Interface Endpoint và AWS
PrivateLink}\par

Đối với những dịch vụ không hỗ trợ Gateway Endpoint, hệ thống có thể sử
dụng Interface Endpoint.

Interface Endpoint tạo các Elastic Network Interface trong subnet và gán
địa chỉ IP nội bộ cho chúng. Ứng dụng có thể sử dụng private IP hoặc
private DNS để truy cập dịch vụ AWS mà không cần đi qua Internet công
cộng.

Một số dịch vụ thường được truy cập thông qua Interface Endpoint gồm:

\begin{itemize}
\tightlist
\item
  Amazon CloudWatch.
\item
  AWS Secrets Manager.
\item
  Amazon ECR.
\item
  AWS Systems Manager.
\item
  AWS Key Management Service.
\item
  Amazon Kinesis.
\end{itemize}

Interface Endpoint sử dụng AWS PrivateLink để duy trì lưu lượng trong
mạng AWS.

So với việc truy cập thông qua NAT Gateway, giải pháp này có thể:

\begin{itemize}
\tightlist
\item
  Giảm lưu lượng Internet outbound.
\item
  Kiểm soát truy cập bằng Security Group.
\item
  Hạn chế phạm vi tấn công.
\item
  Hỗ trợ kiến trúc private networking rõ ràng hơn.
\item
  Cho phép workload trong Private Subnet truy cập dịch vụ AWS bằng
  private IP.
\end{itemize}

Interface Endpoint vẫn có chi phí theo giờ và lưu lượng xử lý, vì vậy
cần so sánh với chi phí NAT Gateway dựa trên lưu lượng thực tế của hệ
thống.

\vspace{0.5em}\noindent\textbf{2. Triển khai Multi-AZ và tránh chi phí Cross-AZ không cần
thiết}\par

\vspace{0.5em}\noindent\textbf{Tại sao cần nhiều Availability
Zones?}\par

Một Availability Zone có thể gặp sự cố độc lập. Nếu toàn bộ application
server, database và network egress đều được đặt trong cùng một AZ, hệ
thống có thể bị gián đoạn khi AZ đó gặp vấn đề.

Để tăng tính sẵn sàng, kiến trúc production thường được triển khai trên
ít nhất hai Availability Zones.

Ví dụ:

\begin{Shaded}
\begin{Highlighting}[]
\NormalTok{Availability Zone A}
\NormalTok{{-} Public Subnet A}
\NormalTok{{-} Private Subnet A}
\NormalTok{{-} NAT Gateway A}
\NormalTok{{-} Application instances A}

\NormalTok{Availability Zone B}
\NormalTok{{-} Public Subnet B}
\NormalTok{{-} Private Subnet B}
\NormalTok{{-} NAT Gateway B}
\NormalTok{{-} Application instances B}
\end{Highlighting}
\end{Shaded}

Trong mô hình này, mỗi Availability Zone có một đường outbound độc lập.
Nếu một AZ gặp sự cố, workload trong AZ còn lại vẫn có thể sử dụng NAT
Gateway và các tài nguyên còn hoạt động trong AZ đó.

\vspace{0.5em}\noindent\textbf{Tránh Single Point of
Failure}\par

Nếu hệ thống chỉ sử dụng một NAT Gateway trong AZ-A, các workload nằm
trong AZ-B phải truyền lưu lượng sang AZ-A trước khi truy cập Internet.

Mô hình này tạo ra hai vấn đề:

\begin{enumerate}
\def\labelenumi{\arabic{enumi}.}
\tightlist
\item
  NAT Gateway trong AZ-A trở thành điểm phụ thuộc duy nhất.
\item
  Lưu lượng từ AZ-B sang AZ-A có thể phát sinh chi phí truyền dữ liệu
  Cross-AZ.
\end{enumerate}

Một kiến trúc phù hợp hơn là triển khai một NAT Gateway trong mỗi
Availability Zone đang hoạt động. Route table của từng Private Subnet sẽ
trỏ đến NAT Gateway nằm trong cùng AZ.

Ví dụ:

\begin{Shaded}
\begin{Highlighting}[]
\NormalTok{Private Subnet A {-}\textgreater{} NAT Gateway A}
\NormalTok{Private Subnet B {-}\textgreater{} NAT Gateway B}
\end{Highlighting}
\end{Shaded}

Nếu AZ-A gặp sự cố, tài nguyên trong AZ-B vẫn có đường outbound riêng và
không phụ thuộc vào NAT Gateway A.

\vspace{0.5em}\noindent\textbf{AZ-Aware Routing}\par

AZ-Aware Routing là cách ưu tiên kết nối giữa các workload và dependency
nằm trong cùng Availability Zone.

Ví dụ:

\begin{itemize}
\tightlist
\item
  Application instance trong AZ-A ưu tiên sử dụng NAT Gateway A.
\item
  Workload trong AZ-B ưu tiên kết nối cache node trong AZ-B.
\item
  Load balancer phân phối request đến các target phù hợp.
\item
  Database replica hoặc cache node được phân bố giữa nhiều AZ.
\item
  Service discovery ưu tiên endpoint gần workload hơn khi kiến trúc hỗ
  trợ.
\end{itemize}

Cách bố trí này giúp:

\begin{itemize}
\tightlist
\item
  Giảm lưu lượng Cross-AZ không cần thiết.
\item
  Giảm độ trễ network.
\item
  Hạn chế chi phí truyền dữ liệu giữa các AZ.
\item
  Giúp mỗi AZ có khả năng vận hành độc lập tốt hơn.
\end{itemize}

Tuy nhiên, việc ưu tiên kết nối trong cùng AZ vẫn phải được cân bằng với
yêu cầu High Availability. Ứng dụng không nên phụ thuộc hoàn toàn vào
một tài nguyên trong cùng AZ nếu tài nguyên đó không có cơ chế failover.

\vspace{0.5em}\noindent\textbf{3. Quản lý kết nối nhiều VPC bằng AWS Transit
Gateway}\par

\vspace{0.5em}\noindent\textbf{Hạn chế của VPC
Peering}\par

VPC Peering phù hợp khi hệ thống chỉ có một số lượng nhỏ VPC. Hai VPC có
thể được kết nối trực tiếp và giao tiếp với nhau bằng private IP.

Tuy nhiên, VPC Peering không hỗ trợ transitive routing.

Ví dụ:

\begin{Shaded}
\begin{Highlighting}[]
\NormalTok{VPC{-}A \textless{}{-}\textgreater{} VPC{-}B}
\NormalTok{VPC{-}B \textless{}{-}\textgreater{} VPC{-}C}
\end{Highlighting}
\end{Shaded}

Việc VPC-A kết nối với VPC-B và VPC-B kết nối với VPC-C không có nghĩa
là VPC-A có thể tự động giao tiếp với VPC-C.

Nếu nhiều VPC cần kết nối trực tiếp với nhau, số lượng VPC Peering
connections có thể tăng theo công thức:

\begin{Shaded}
\begin{Highlighting}[]
\NormalTok{N(N {-} 1) / 2}
\end{Highlighting}
\end{Shaded}

Ví dụ, với 10 VPC, mô hình full mesh có thể cần đến 45 kết nối peering.

Khi đó, việc quản lý route table, CIDR, Security Group và network policy
trở nên phức tạp. Việc thêm một VPC mới cũng có thể yêu cầu tạo nhiều
kết nối và cập nhật nhiều route table khác nhau.

\vspace{0.5em}\noindent\textbf{Transit Gateway như một Cloud
Router}\par

AWS Transit Gateway hoạt động như một router trung tâm cho nhiều VPC và
hệ thống mạng khác nhau.

Các thành phần có thể kết nối với Transit Gateway gồm:

\begin{itemize}
\tightlist
\item
  Amazon VPC.
\item
  Site-to-Site VPN.
\item
  AWS Direct Connect thông qua Direct Connect Gateway.
\item
  VPC thuộc nhiều AWS accounts.
\item
  Các môi trường Development, Staging và Production.
\item
  Shared Services VPC.
\item
  Mạng từ Data Center On-Premises.
\end{itemize}

Thay vì tạo kết nối trực tiếp giữa từng cặp VPC, mỗi VPC chỉ cần tạo một
attachment đến Transit Gateway.

\begin{Shaded}
\begin{Highlighting}[]
\NormalTok{             Development VPC}
\NormalTok{                    |}
\NormalTok{                    |}
\NormalTok{Shared VPC {-}{-}{-}{-} Transit Gateway {-}{-}{-}{-} Production VPC}
\NormalTok{                    |}
\NormalTok{                    |}
\NormalTok{             On{-}Premises Network}
\end{Highlighting}
\end{Shaded}

Mô hình hub-and-spoke này giúp đơn giản hóa network topology và tập
trung việc quản lý routing.

\vspace{0.5em}\noindent\textbf{Phân vùng mạng bằng Transit Gateway Route
Tables}\par

Transit Gateway có thể sử dụng nhiều route table để kiểm soát VPC nào
được phép giao tiếp với VPC nào.

Ví dụ:

\begin{itemize}
\tightlist
\item
  Development VPC có thể truy cập Shared Services VPC.
\item
  Production VPC có thể truy cập Shared Services VPC.
\item
  Development VPC không được phép kết nối trực tiếp đến Production VPC.
\item
  On-Premises Network chỉ được truy cập một số subnet nhất định.
\item
  Security VPC có thể kiểm tra lưu lượng trước khi chuyển đến các VPC
  khác.
\end{itemize}

Việc phân vùng này giúp:

\begin{itemize}
\tightlist
\item
  Hạn chế lateral movement giữa các môi trường.
\item
  Phân tách Development, Staging và Production rõ ràng hơn.
\item
  Quản lý network policy từ một vị trí tập trung.
\item
  Đơn giản hóa việc thêm VPC mới.
\item
  Hỗ trợ kiến trúc multi-account với AWS Organizations.
\end{itemize}

\vspace{0.5em}\noindent\textbf{So sánh mô hình VPC cơ bản và kiến trúc nâng
cao}\par

{
\begingroup
\small
\setlength{\tabcolsep}{3pt}
\renewcommand{\arraystretch}{1.15}
\begin{longtable}{@{}
  >{\raggedright\arraybackslash}p{(\linewidth - 4\tabcolsep) * \real{0.3333}}
  >{\raggedright\arraybackslash}p{(\linewidth - 4\tabcolsep) * \real{0.3333}}
  >{\raggedright\arraybackslash}p{(\linewidth - 4\tabcolsep) * \real{0.3333}}@{}}
\toprule\noalign{}
\begin{minipage}[b]{\linewidth}\raggedright
Tiêu chí
\end{minipage} & \begin{minipage}[b]{\linewidth}\raggedright
Mô hình cơ bản
\end{minipage} & \begin{minipage}[b]{\linewidth}\raggedright
Mô hình nâng cao
\end{minipage} \\
\midrule\noalign{}
\endhead
\bottomrule\noalign{}
\endlastfoot
Truy cập S3 và DynamoDB & Đi qua NAT Gateway hoặc public endpoint & Đi
qua Gateway Endpoint \\
Truy cập các AWS services khác & Thường sử dụng NAT Gateway & Sử dụng
Interface Endpoint và AWS PrivateLink khi phù hợp \\
High Availability & Có thể chỉ sử dụng một NAT Gateway & NAT Gateway độc
lập trong từng Availability Zone \\
Cross-AZ traffic & Dễ phát sinh nếu routing không được tối ưu & Hạn chế
bằng AZ-aware routing \\
Kết nối nhiều VPC & VPC Peering được cấu hình thủ công từng cặp &
Transit Gateway quản lý tập trung \\
Phân tách Dev và Prod & Phụ thuộc vào nhiều route và peering connections
& Sử dụng Transit Gateway route tables \\
Khả năng mở rộng & Phù hợp với hệ thống nhỏ & Phù hợp hơn với hệ thống
nhiều VPC và nhiều account \\
Quản trị & Dễ triển khai ban đầu nhưng khó kiểm soát khi mở rộng & Phức
tạp hơn lúc thiết kế nhưng dễ quản lý tập trung \\
\end{longtable}
\endgroup
}

\vspace{0.5em}\noindent\textbf{Khi nào nên sử dụng từng giải
pháp?}\par

\vspace{0.5em}\noindent\textbf{Sử dụng VPC Endpoints
khi:}\par

\begin{itemize}
\tightlist
\item
  Workload trong Private Subnet thường xuyên truy cập Amazon S3 hoặc
  DynamoDB.
\item
  Ứng dụng cần truy cập AWS services mà không đi qua Internet công cộng.
\item
  Muốn giảm lưu lượng và chi phí xử lý qua NAT Gateway.
\item
  Hệ thống có yêu cầu cao về bảo mật dữ liệu.
\item
  Muốn giới hạn quyền truy cập bằng endpoint policy.
\end{itemize}

\vspace{0.5em}\noindent\textbf{Sử dụng Multi-AZ NAT Gateway
khi:}\par

\begin{itemize}
\tightlist
\item
  Hệ thống yêu cầu High Availability.
\item
  Workload được triển khai trên nhiều Availability Zones.
\item
  Không muốn một NAT Gateway trở thành Single Point of Failure.
\item
  Muốn hạn chế lưu lượng Cross-AZ không cần thiết.
\item
  Hệ thống cần duy trì outbound connectivity khi một AZ gặp sự cố.
\end{itemize}

\vspace{0.5em}\noindent\textbf{Sử dụng AWS Transit Gateway
khi:}\par

\begin{itemize}
\tightlist
\item
  Hệ thống có nhiều VPC hoặc nhiều AWS accounts.
\item
  Cần kết nối mạng On-Premises với AWS.
\item
  VPC Peering bắt đầu trở nên khó quản lý.
\item
  Cần phân tách Development, Staging và Production.
\item
  Muốn kiểm soát routing từ một vị trí trung tâm.
\item
  Cần xây dựng mô hình Shared Services hoặc Security Inspection tập
  trung.
\end{itemize}

\vspace{0.5em}\noindent\textbf{Lưu ý về chi phí}\par

Kiến trúc network nâng cao không có nghĩa là mọi thành phần đều rẻ hơn.

Gateway Endpoint cho S3 và DynamoDB có thể giúp giảm lưu lượng qua NAT
Gateway. Tuy nhiên, Interface Endpoint, NAT Gateway và Transit Gateway
đều có mô hình tính phí riêng.

Trước khi triển khai, cần phân tích:

\begin{itemize}
\tightlist
\item
  Tổng lưu lượng dữ liệu hàng tháng.
\item
  Các dịch vụ AWS được truy cập thường xuyên.
\item
  Số lượng Availability Zones.
\item
  Số lượng VPC và Transit Gateway attachments.
\item
  Tần suất truyền dữ liệu giữa các AZ.
\item
  Lưu lượng Internet outbound.
\item
  Yêu cầu bảo mật và tính sẵn sàng.
\item
  Chi phí vận hành và độ phức tạp trong quản trị.
\end{itemize}

Mục tiêu không phải là giảm chi phí bằng mọi cách, mà là tìm được sự cân
bằng giữa chi phí, bảo mật, hiệu năng, khả năng mở rộng và High
Availability.

\vspace{0.5em}\noindent\textbf{Kết luận}\par

Mô hình Public Subnet, Private Subnet và NAT Gateway là nền tảng tốt để
bắt đầu xây dựng hệ thống trên AWS. Tuy nhiên, khi hệ thống phát triển
thành nhiều môi trường, nhiều VPC hoặc nhiều AWS accounts, network
architecture cũng cần được nâng cấp.

Ba kỹ thuật quan trọng gồm:

\begin{enumerate}
\def\labelenumi{\arabic{enumi}.}
\tightlist
\item
  Sử dụng VPC Endpoints để giữ lưu lượng truy cập AWS services trong
  mạng nội bộ và giảm phụ thuộc vào NAT Gateway.
\item
  Triển khai Multi-AZ với NAT Gateway độc lập trong từng Availability
  Zone để tăng tính sẵn sàng và hạn chế Cross-AZ traffic.
\item
  Sử dụng AWS Transit Gateway để quản lý kết nối giữa nhiều VPC theo mô
  hình tập trung.
\end{enumerate}

Tối ưu kiến trúc VPC không chỉ là bài toán routing. Đây còn là quá trình
cân bằng giữa bảo mật, khả năng mở rộng, tính sẵn sàng, hiệu năng và chi
phí.

Một Enterprise VPC Architecture tốt cần được thiết kế dựa trên nhu cầu
thực tế của hệ thống. Không nên bổ sung thêm nhiều dịch vụ chỉ để kiến
trúc trở nên phức tạp hơn mà không tạo ra giá trị rõ ràng.

\vspace{0.5em}\noindent\textbf{Thông tin bài đăng}\par

\begin{itemize}
\tightlist
\item
  \textbf{Chủ đề:} Nâng cấp kiến trúc Amazon VPC cho hệ thống doanh
  nghiệp.
\item
  \textbf{Ngày đăng:} 29/07/2026.
\item
  \textbf{Nền tảng:} AWS Study Groups.
\item
  \textbf{Trạng thái:} Published.
\item
  \textbf{Public link:}
  \href{https://www.facebook.com/groups/awsstudygroupfcj/permalink/2225997548165205/\#}{Blog
  4 trên AWS Study Groups}
\end{itemize}
